\section{Results}\label{sec:results}

\subsection{Seasonal mismatch}\label{sec:monthly}
In 2024 Italian solar plants produced 1{,}541\,GWh in January and 4{,}700\,GWh in July; wind produced 2{,}814\,GWh in January and 916\,GWh in August \citep{GSE2025,Terna2025m} (Fig.~\ref{fig:profiles}). Monthly costs in this subsection are per MWh of total demand (Section~\ref{sec:methods}). Covering the worst month with solar alone requires 1.95 times the capacity sized on the annual mean, so the real lifetime cost of 42\,EUR/MWh derived from the main FER~X auction becomes 82\,EUR/MWh of useful energy before any storage (96 at the NZIA price). Adding daily batteries at the MACSE price gives 110\,EUR/MWh (125). Covering the seasonal deficit, 16.8\% of annual energy, with batteries that cycle once a year costs about 2{,}000\,EUR/MWh; with hydrogen stored in salt caverns at the cost estimated for purely seasonal use \citep{DCU2026}, about 390\,EUR/MWh, a lower bound because the cost of electrolysers and fuel cells is excluded. A 50/50 mix of solar and wind reduces the overbuild to 1.35 and the seasonal deficit to 5\%, for 76--92\,EUR/MWh with daily batteries (84 at the reference wind price), or 90--106 with a grid cost of 14\,EUR/MWh. On this monthly view, without grid costs, the mix is cheaper than nuclear at the Dukovany price financed at 5\% (<<LC:8540:5:7|0>>\,EUR/MWh at a 90\% load factor, Table~\ref{tab:lcoe}). The monthly figures, however, ignore the hours and weeks without wind and sun \citep{KittelSchill2024} and heat waves, during which wind output in southern Europe can fall by up to 30.8\% while demand rises \citep{Molina2023}; a monitoring of large solar parks by the University of \'Evora, reported by the press, measured an average loss of 17.5\% of output over 24 hours in extreme heat, because module efficiency falls above 25\,$^\circ$C \citep{Evora2026}. The three-hour model includes them for the weather years considered, and the comparison changes.

\begin{figure}[t]
\centering\includegraphics[width=\linewidth]{fig1_profiles.pdf}
\caption{Monthly output of solar PV and wind in Italy in 2024, as a share of each technology's annual output, against flat demand. Data: \citet{GSE2025}, \citet{Terna2025m}.}\label{fig:profiles}
\end{figure}

\subsection{One curve: from fixed cost to nuclear capacity}\label{sec:least}
The central result is one curve (Fig.~\ref{fig:share}): the annualised fixed cost of nuclear, fixN in EUR/kW-yr, against the nuclear capacity of the least-cost system under net self-sufficiency (Section~\ref{sec:model}), at a 5\% and a 7\% cost of capital, for the PI and OB cost cases; the curve of the earlier import rule (imports with a capacity credit and no exports) is shown as a variant. Cost cases are points on the curve, not separate runs: each is converted to fixN with the formula of Section~\ref{sec:model} and read on the curve (Table~\ref{tab:cases}). The curve points are six-hour runs; the three-hour central points are shown as separate markers. Costs are per MWh of demand not met by existing hydropower.

\begin{figure}[t]
\centering\includegraphics[width=0.95\linewidth]{fig2_share_vs_fixedcost.pdf}
\caption{Optimal nuclear capacity against the annualised fixed cost of nuclear (fixN) under net self-sufficiency, at 5\% (blue) and 7\% (orange). Solid lines with circles: first plant (PI); dashed with squares: target (OB, reduced model). Thin grey lines: the earlier rule, imports with a capacity credit, at 5\%. Vertical lines: the five cost cases at 5\% and the central case at 7\%. Lines are six-hour runs; diamonds and crosses are the three-hour central points. Source: \texttt{results/v7.json} and \texttt{v7\_sweep.json}.}\label{fig:share}
\end{figure}

\begin{table}[htbp]
\centering
\caption{The cost cases as points on the curve, six-hour step: optimal nuclear capacity in GW and, in parentheses, share of total demand in \%, under net self-sufficiency and, as a variant, with imports with a capacity credit. fixN is 669.5 at 5\% and 922.7 at 7\% for Dukovany (runs at 670 and 923), 878.3 at 5\% and 1{,}334.2 at 7\% for the risk case. * interpolated between two computed points, or 0 beyond the last point with nuclear.}\label{tab:cases}
\input{tab_cases}
\end{table}

Under net self-sufficiency at 5\%, nuclear falls from <<N6:PI:s:5:335|1>>\,GW at Saeul cost (fixN 335, PI) to <<N6:PI:s:5:598|1>>\,GW at 598, to <<N6:PI:s:5:604|1>>\,GW at the EPC-only price (604), to <<RUN:ss0_PI_670_h6:N_GW|1>>\,GW in the central case (670) and to <<N6:PI:s:5:750|1>>\,GW at 750; a tail between 12 and 9\,GW follows up to about 880, and nuclear disappears between 950 and 1{,}050. The slope is steepest between 600 and 750 (about <<EXPR:(V('N6:PI:s:5:598')-V('N6:PI:s:5:750'))/152|2>>\,GW per EUR/kW-yr, against <<EXPR:(V('N6:PI:s:5:335')-V('N6:PI:s:5:598'))/263|2>> between 335 and 598): this is where the central case lies, and a 10\% error on construction cost moves the result by several GW. The OB curve falls earlier: <<N6:OB:s:5:750|1>>\,GW at 750 against <<N6:PI:s:5:750|1>> in PI, and zero from 950, because in OB the technologies at target cost lower the cost of the system without nuclear (<<RUN:ss0_OB_none_h6:cost_new|2>> against <<RUN:ss0_PI_none_h6:cost_new|2>>\,EUR/MWh). At 670 the two series give <<RUN:ss0_PI_670_h6:N_GW|1>> and <<RUN:ss0_OB_670_h6:N_GW|1>>\,GW. The risk case (9{,}680\,EUR/kW, 15 years, fixN 878.3) gives <<N6:PI:s:5:878|1>>\,GW in PI and <<N6:OB:s:5:878|1>> in OB; Vogtle gives none. Fifteen years of construction alone, at the central cost (fixN <<FIX:8540:5:15|1>>), give <<N6:PI:s:5:792.5|1>>\,GW in PI and <<N6:OB:s:5:792.5|1>>\,GW in OB at six hours, interpolated. Barakah (4{,}670\,EUR/kW, 8 years, fixN <<FIX:4670:5:8|1>>) gives <<N6:PI:s:5:442|1>>\,GW in PI and <<N6:OB:s:5:442|1>> in OB. Under the earlier import rule the curve is lower in the central region: <<RUN:c3_PI_670_h6:N_GW|1>> and <<RUN:c3_OB_670_h6:N_GW|1>>\,GW at 670.

\subsection{Thresholds}
Table~\ref{tab:thr} gives the fixN at which nuclear falls below 10\,GW, by linear interpolation between the two computed points that bracket 10\,GW, and the equivalent construction cost with 7 and 15 years of construction. Under net self-sufficiency at 5\% and a six-hour step the PI threshold is <<THR6:PI:s:5:10|0>> (<<EXPR:V('CPX:'+str(V('THR6:PI:s:5:10'))+':5:7')|0>>\,EUR/kW at 7 years) and the OB threshold <<THR6:OB:s:5:10|0>> (<<EXPR:V('CPX:'+str(V('THR6:OB:s:5:10'))+':5:7')|0>>). A construction cost of about 10{,}000\,EUR/kW, the range of Flamanville, keeps PI above 10\,GW and brings OB below. The central case lies <<EXPR:V('THR6:PI:s:5:10')-670|0>>\,EUR/kW-yr below the PI threshold, about <<EXPR:V('CPX:'+str(V('THR6:PI:s:5:10'))+':5:7')-8540|-2>>\,EUR/kW of construction cost, and <<EXPR:V('THR6:OB:s:5:10')-670|0>> below the OB one. At three hours the PI threshold rises: PI still has <<RUN:ss0_PI_850_h3:N_GW|1>>\,GW at fixN 850 (11{,}500\,EUR/kW at seven years), so the threshold lies beyond 850, against <<THR6:PI:s:5:10|0>> at six hours. For OB the three-hour points do not pin it down: OB has <<RUN:ss0_OB_670_h3:N_GW|1>>\,GW at 670 and <<RUN:ss0_OB_780_h3:N_GW|1>> at 780 (<<N6:OB:s:5:780|1>> at six hours), so the threshold lies between 670 and 780 as at six hours, and it cannot be said to rise. Between fixN 780 and 850 three-hour nuclear barely changes, from <<RUN:ss0_PI_780_h3:N_GW|1>> to <<RUN:ss0_PI_850_h3:N_GW|1>>\,GW in PI and from <<RUN:ss0_OB_780_h3:N_GW|1>> to <<RUN:ss0_OB_850_h3:N_GW|1>> in OB (Table~\ref{tab:step}). With 15 years of construction the thresholds fall to about <<EXPR:V('CPX:'+str(V('THR6:PI:s:5:10'))+':5:15')|-2>> and <<EXPR:V('CPX:'+str(V('THR6:OB:s:5:10'))+':5:15')|-2>>\,EUR/kW. At 7\% the PI threshold is fixN <<THR6:PI:s:7:10|0>> (<<EXPR:V('CPX:'+str(V('THR6:PI:s:7:10'))+':7:7')|0>>\,EUR/kW at 7 years): the central case, 8{,}540, lies just below it and has <<RUN:ss0_PI_7pc_923_h3:N_GW|1>>\,GW in PI and <<RUN:ss0_OB_7pc_923_h3:N_GW|1>> in OB at three hours.

\begin{table}[htbp]
\centering
\caption{fixN at which optimal nuclear capacity falls below 10\,GW, by import rule, and equivalent construction cost with 7 and 15 years of construction. * three-hour OB threshold under the capacity-credit rule: linear extrapolation, a lower bound.}\label{tab:thr}
\input{tab_thr}
\end{table}

\subsection{The time step and the central values}\label{sec:step}
Going from a six-hour to a three-hour step raises the cost of the system without nuclear (by <<EXPR:V('RUN:ss0_PI_none_h3:cost_new')-V('RUN:ss0_PI_none_h6:cost_new')|1>>\,EUR/MWh in PI) more than that of the system with nuclear (<<EXPR:V('RUN:ss0_PI_670_h3:cost_new')-V('RUN:ss0_PI_670_h6:cost_new')|1>>), so the gap grows and, at equal fixN, so does optimal nuclear capacity (Table~\ref{tab:step}). A coarse step averages out short-term variability that solar, wind and batteries must cover and nuclear baseload need not follow; this explanation is plausible but not verified. The six-hour step therefore underestimates the value of nuclear and its optimal capacity at every point computed at three hours (the threshold rises for certain only in PI), and absolute six-hour cost levels are underestimated by an amount that cannot be extrapolated: at two years (run only at fixN 598) the cost of PI rises by 1.60\,EUR/MWh from 12 to 6 hours and by 2.49 from 6 to 3 hours. The gap converges better, and is the decision metric.

\begin{table}[htbp]
\centering
\caption{Costs and capacities at six- and three-hour steps, four years, net self-sufficiency, 5\% unless stated; cost in EUR/MWh.}\label{tab:step}
\input{tab_step}
\end{table}

Central values are the three-hour results (Table~\ref{tab:central}): PI <<RUN:ss0_PI_670_h3:N_GW|1>>\,GW (<<RUN:ss0_PI_670_h3:sh|1>>\% of demand) and OB <<RUN:ss0_OB_670_h3:N_GW|1>>\,GW (<<RUN:ss0_OB_670_h3:sh|1>>\%), with gaps of <<GAP:ss0_PI_none_h3:ss0_PI_670_h3|2>> and <<GAP:ss0_OB_none_h3:ss0_OB_670_h3|2>>\,EUR/MWh (<<GAP:ss0_PI_none_h6:ss0_PI_670_h6|2>> and <<GAP:ss0_OB_none_h6:ss0_OB_670_h6|2>> at six hours); about <<EXPR:V('GAP:ss0_PI_none_h3:ss0_PI_670_h3')*356.6/1000|1>> billion euros a year for PI (<<GAP:ss0_PI_none_h3:ss0_PI_670_h3|2>> times 356.6\,TWh, our calculation). The Richardson extrapolation (Section~\ref{sec:stepmethods}) at fixN 670 with four years has an increment ratio of <<RICHS:PI:670:ratio|2>> for PI, within the acceptable range of 0.35--0.65, and gives a one-hour gap of <<RICHS:PI:670:G1|2>>\,EUR/MWh against <<RICHS:PI:670:g3|2>> at three hours; for OB the ratio is <<RICHS:OB:670:ratio|2>>; under the earlier import rule the PI ratio was <<RICH:PI:670:ratio|2>>, with a one-hour gap of <<RICH:PI:670:G1|2>>\,EUR/MWh against <<RICH:PI:670:g3|2>> at three hours. Where the ratio is outside the range the extrapolation is not reliable: the sign of the correction is certain, its size is not, and the three-hour gap is a lower bound (Table~\ref{tab:conv}). The headline figures use the three-hour values.

\begin{table}[htbp]
\centering
\caption{Central values under net self-sufficiency, and the variant with imports with a capacity credit (earlier rule). Cost in EUR/MWh; share of total demand in \%; gap = cost without nuclear minus cost with nuclear. The central value is the three-hour result.}\label{tab:central}
\input{tab_central}
\end{table}

\begin{table}[htbp]
\centering
\caption{Convergence of the gap with the time step (EUR/MWh). The limit for $h\to0$ is $G_0$ of the Richardson estimate. * ratio outside 0.35--0.65: extrapolation not reliable (n.a., not computed). Rows under the capacity-credit rule at fixN 598 (two and four years) and 670 (four years), and under net self-sufficiency at 670 (four years).}\label{tab:conv}
\input{tab_conv}
\end{table}

\subsection{The price of import dependence, and 7\%}\label{sec:dep}
Table~\ref{tab:dep} compares import rules at six hours, with nuclear at the central cost and without nuclear. The relevant comparison is not between today's import share and that of 2040, but between a system that does not rely on others and one that accepts a fixed share of their energy. With nuclear, allowing net imports of up to 10\% of demand lowers the system cost from <<RUN:ss0_PI_670_h6:cost_new|2>> to <<RUN:ss10_PI_670_h6:cost_new|2>>\,EUR/MWh in PI and from <<RUN:ss0_OB_670_h6:cost_new|2>> to <<RUN:ss10_OB_670_h6:cost_new|2>> in OB, by <<EXPR:V('RUN:ss0_PI_670_h6:cost_new')-V('RUN:ss10_PI_670_h6:cost_new')|2>> and <<EXPR:V('RUN:ss0_OB_670_h6:cost_new')-V('RUN:ss10_OB_670_h6:cost_new')|2>>\,EUR/MWh, and the constraint does not even bind: the system imports <<EXPR:V('UN:ss10_PI_670_h6:IMPORT:TWh')-V('UN:ss10_PI_670_h6:EXPORT:TWh')|0>>\,TWh net in PI out of 40 allowed. Without nuclear the same allowance is worth <<EXPR:V('RUN:ss0_PI_none_h6:cost_new')-V('RUN:ss10_PI_none_h6:cost_new')|2>> (PI) and <<EXPR:V('RUN:ss0_OB_none_h6:cost_new')-V('RUN:ss10_OB_none_h6:cost_new')|2>>\,EUR/MWh (OB): it is the system without nuclear that needs other countries' energy. A 10\% net import target thus buys about 1--1.5\% of system cost with a structural dependence. The earlier rule is cheaper still, <<RUN:c3_PI_670_h6:cost_new|2>>\,EUR/MWh in PI, because it relies on about 7\,GW of foreign capacity in the critical weeks (38\% of the import capacity) and never returns the energy it takes; that value comes from capacity that other countries build and pay for, and would not exist if every country planned the same way.

\begin{table}[htbp]
\centering
\caption{The price of import dependence: system cost at six hours, four years, PI and OB, with nuclear at fixN 670 and without nuclear, by import rule. $X$ is the allowed annual net import as a share of demand. Net imports are imports minus exports in the case with nuclear. Cost in EUR/MWh.}\label{tab:dep}
\input{tab_dep}
\end{table}

At a 7\% cost of capital the central case (fixN 923) gives <<RUN:ss0_PI_7pc_923_h6:N_GW|1>>\,GW in PI and <<RUN:ss0_OB_7pc_923_h6:N_GW|1>> in OB at six hours, <<RUN:ss0_PI_7pc_923_h3:N_GW|1>> and <<RUN:ss0_OB_7pc_923_h3:N_GW|1>> at three hours, with gaps of <<GAP:ss0_PI_7pc_none_h3:ss0_PI_7pc_923_h3|2>> and <<GAP:ss0_OB_7pc_none_h3:ss0_OB_7pc_923_h3|2>>\,EUR/MWh at three hours. At 7\% the costs of solar, wind, storage and hydrogen also rise, so at equal fixN the 7\% curve lies above the 5\% one; it is the fixN of nuclear that rises most, from 670 to 923. The risk case at 7\% (fixN 1{,}334.2) gives no nuclear; nuclear is zero from 1{,}100. Without imports of any kind, a stricter variant of self-sufficiency, the central case gives <<N6:PI:n:5:670|1>>\,GW in PI and <<N6:OB:n:5:670|1>> in OB at 5\%, and about <<N6:PI:n:7:923|1>> and <<N6:OB:n:7:923|1>> at 7\% (six hours, interpolated). The full points are in Tables~\ref{tab:curve5} and~\ref{tab:other}.

\subsection{The central system}
At three hours the central PI system (fixN 670) has <<UN:ss0_PI_670_h3:NUC:GW|1>>\,GW of nuclear (<<UN:ss0_PI_670_h3:NUC:hours|0>> equivalent hours, <<UN:ss0_PI_670_h3:NUC:TWh|1>>\,TWh), <<GR:ss0_PI_670_h3:pv:GW|1>>\,GW of solar, <<GR:ss0_PI_670_h3:on:GW|1>>\,GW of onshore wind, <<UN:ss0_PI_670_h3:IMPORT:GW|1>>\,GW of interconnection (<<UN:ss0_PI_670_h3:IMPORT:TWh|1>>\,TWh imported and as much exported) and <<GR:ss0_PI_670_h3:battery:GWh|0>>\,GWh of batteries. At six hours it has <<UN:ss0_PI_670_h6:NUC:GW|1>>\,GW of nuclear (<<UN:ss0_PI_670_h6:NUC:hours|0>>\,h, <<UN:ss0_PI_670_h6:NUC:TWh|1>>\,TWh), <<GR:ss0_PI_670_h6:pv:GW|1>>\,GW of solar (24.2 existing, <<EXPR:V('GR:ss0_PI_670_h6:pv:GW')-V('UN:ss0_PI_670_h6:PVex_terra:GW')-V('UN:ss0_PI_670_h6:PVex_tetti:GW')|1>> new: <<EXPR:V('UN:ss0_PI_670_h6:PV_terra_Sud-isole:GW')+V('UN:ss0_PI_670_h6:PV_terra-pipeline_Lazio/Sicilia/Piem./Puglia/Sard.:GW')+V('UN:ss0_PI_670_h6:PV_terra_Centro-Sud:GW')+V('UN:ss0_PI_670_h6:PV_terra_Nord:GW')|1>> ground-mounted, all ground blocks exhausted and <<UN:ss0_PI_670_h6:PV_agrivoltaico_Sud-isole:GW|1>> agrivoltaic), <<GR:ss0_PI_670_h6:on:GW|1>>\,GW of onshore wind (E0 13.6, E1 6.0, E2 <<UN:ss0_PI_670_h6:E2:GW|1>>), <<UN:ss0_PI_670_h6:IMPORT:GW|1>>\,GW of interconnection (<<UN:ss0_PI_670_h6:IMPORT:TWh|1>>\,TWh imported and as much exported) and <<GR:ss0_PI_670_h6:battery:GWh|0>>\,GWh of batteries; solar and wind curtailment is about <<EXPR:(FULL['ss0_PI_670_h6']['diag']['VRE']['on']['curt']+FULL['ss0_PI_670_h6']['diag']['VRE']['pv']['curt'])/(FULL['ss0_PI_670_h6']['diag']['VRE']['on']['avail']+FULL['ss0_PI_670_h6']['diag']['VRE']['pv']['avail'])*100|1>>\% of available output. Its cost of <<RUN:ss0_PI_670_h6:cost_new|2>>\,EUR/MWh comprises nuclear <<COMP:ss0_PI_670_h6:nuclear|2>>, solar <<COMP:ss0_PI_670_h6:pv|2>>, wind <<COMP:ss0_PI_670_h6:on|2>>, imports <<COMP:ss0_PI_670_h6:import|2>>, batteries <<COMP:ss0_PI_670_h6:battery|2>>, geothermal <<COMP:ss0_PI_670_h6:geo|2>>, bioenergy <<COMP:ss0_PI_670_h6:bio|2>>, 100-hour storage <<COMP:ss0_PI_670_h6:ldes|2>> and interruptible load <<COMP:ss0_PI_670_h6:interr|2>>. The central OB system at six hours has <<UN:ss0_OB_670_h6:NUC:GW|1>>\,GW of nuclear, <<GR:ss0_OB_670_h6:pv:GW|1>>\,GW of solar, <<GR:ss0_OB_670_h6:on:GW|1>>\,GW of onshore wind, <<GR:ss0_OB_670_h6:battery:GWh|0>>\,GWh of batteries and <<GR:ss0_OB_670_h6:caes:GWh|0>>\,GWh of compressed air (Fig.~\ref{fig:capacity}).

\begin{figure}[t]
\centering\includegraphics[width=0.95\linewidth]{fig3_capacity.pdf}
\caption{Installed capacity and cost of the least-cost system at a six-hour step for PI and OB, with nuclear at fixN 670 (5\%), at 923 (7\%) and without nuclear, under net self-sufficiency. PV includes the existing fleet after 2016; onshore wind includes E0.}\label{fig:capacity}
\end{figure}

\begin{figure}[t]
\centering\includegraphics[width=0.7\linewidth]{fig4_cost_vs_nuclear.pdf}
\caption{Gap (cost without nuclear minus cost with nuclear) against fixN under net self-sufficiency, at 5\% and 7\%. Lines are six-hour runs; isolated points are the three-hour central values (PI <<GAP:ss0_PI_none_h3:ss0_PI_670_h3|2>> and OB <<GAP:ss0_OB_none_h3:ss0_OB_670_h3|2>>\,EUR/MWh at 670; PI <<GAP:ss0_PI_7pc_none_h3:ss0_PI_7pc_923_h3|2>> and OB <<GAP:ss0_OB_7pc_none_h3:ss0_OB_7pc_923_h3|2>> at 7\%, fixN 923).}\label{fig:gap}
\end{figure}

\subsection{The system without nuclear and the bridge from the previous version}\label{sec:bridge}
Under net self-sufficiency, at six hours the PI system without nuclear has <<GR:ss0_PI_none_h6:pv:GW|0>>\,GW of solar, <<GR:ss0_PI_none_h6:on:GW|1>>\,GW of onshore wind, <<UN:ss0_PI_none_h6:IMPORT:GW|1>>\,GW of interconnection (<<UN:ss0_PI_none_h6:IMPORT:TWh|1>>\,TWh imported and as much exported), <<GR:ss0_PI_none_h6:battery:GWh|0>>\,GWh of batteries, <<GR:ss0_PI_none_h6:h2:GWh|0>>\,GWh of hydrogen storage with a <<UN:ss0_PI_none_h6:H2:Pd_GW|1>>\,GW turbine and <<GR:ss0_PI_none_h6:ldes:GWh|0>>\,GWh of 100-hour storage; OB has <<GR:ss0_OB_none_h6:pv:GW|0>>\,GW of solar, <<GR:ss0_OB_none_h6:on:GW|1>>\,GW of onshore wind, <<GR:ss0_OB_none_h6:caes:GWh|0>>\,GWh of compressed air, <<GR:ss0_OB_none_h6:h2:GWh|0>>\,GWh of hydrogen storage and <<GR:ss0_OB_none_h6:battery:GWh|0>>\,GWh of batteries. It costs <<RUN:ss0_PI_none_h6:cost_new|2>> (PI) and <<RUN:ss0_OB_none_h6:cost_new|2>>\,EUR/MWh (OB) at six hours, <<RUN:ss0_PI_none_h3:cost_new|2>> and <<RUN:ss0_OB_none_h3:cost_new|2>> at three. The previous version of this paper, whose model lacked these technologies, gave 204.2\,EUR/MWh and 38.5\,GW of nuclear supplying 59\% of demand. Table~\ref{tab:bridgepi} rebuilds the difference one technology at a time at six hours and with nuclear at fixN 598, the cost of the previous version, which is not the present central case (670); step 0 reproduces the previous version in the present code (193.01\,EUR/MWh without nuclear and 94.05 with nuclear at fixN 598, 38.4\,GW; 204.2 is the same model at three hours).

\begin{table}[htbp]
\centering
\caption{From the previous model to the present one, PI case, six hours, cumulative; cost in EUR/MWh. The last column is the change in cost of the system with nuclear. The OB case is Table~\ref{tab:bridgeob}.}\label{tab:bridgepi}
\input{tab_bridge_pi}
\end{table}

Imports explain most of the reduction without nuclear: 72.67 of 103.20\,EUR/MWh (193.01 minus 89.81). Hydrogen with electrolysis, geological storage and turbine removes 16.00, bioenergy 7.81, geothermal 5.11, smart charging and interruptible load 0.59 and 100-hour storage 1.00; concentrated solar power, compressed air and new pumped hydro are worth nothing in PI, while in OB the first is worth 1.91 and the second 3.80. With nuclear at fixN 598 (as in the table), nuclear falls from 38.4\,GW to 27.9 with imports, 26.4 with geothermal and 23.3 with bioenergy, and rises to 23.9 with smart charging; hydrogen does not change optimal nuclear. The bridge uses the earlier import rule (capacity credit, no exports); the last step, from that rule to net self-sufficiency, raises the cost without nuclear from 89.81 to <<RUN:ss0_PI_none_h6:cost_new|2>>\,EUR/MWh and that with nuclear at fixN 598 from <<RUN:sw_PI_imp_598:cost_new|2>> to <<RUN:ss0_PI_598_h6:cost_new|2>>, and nuclear from <<RUN:sw_PI_imp_598:N_GW|1>> to <<RUN:ss0_PI_598_h6:N_GW|1>>\,GW (PI, six hours). Large-scale geological hydrogen storage exists only as studies and pilots, and the hydrogen-to-power chain has no commercial plant: model costs are first-plant or target values from studies, not observed prices. Without nuclear hydrogen is worth 16\,EUR/MWh, the most uncertain part of the cost of the system without nuclear; the gap is computed with those costs.

\subsection{Sensitivity analyses}\label{sec:sens}
The sensitivities of this subsection were computed under the earlier import rule (capacity credit, no exports) and were not repeated under net self-sufficiency. They are compared with the central case of the same rule at six hours: <<RUN:c3_PI_670_h6:N_GW|1>>\,GW and a gap of <<GAP:val_PI_none:c3_PI_670_h6|2>>\,EUR/MWh in PI, <<RUN:c3_OB_670_h6:N_GW|1>>\,GW and <<GAP:val_OB_none:c3_OB_670_h6|2>> in OB. Under net self-sufficiency the base gap is higher (<<GAP:ss0_PI_none_h6:ss0_PI_670_h6|2>>\,EUR/MWh in PI and <<GAP:ss0_OB_none_h6:ss0_OB_670_h6|2>> in OB), but the effect of each variation under the new rule was not recomputed and its size is unknown: CSP, gas with DAC or SMRs could matter more, because they would also replace the balancing that imports provided before.

\paragraph{Concentrated solar power.} CSP enters only in OB and, in the base case, only without nuclear; its 10\,GW ceiling is our assumption. With nuclear at fixN 670 the OB gap is <<GAP:val_OB_none:c3_OB_670_h6|2>>\,EUR/MWh; a 3\,GW ceiling gives <<GAP:se_OB_csp_c3_none:c3s_se_OB_csp_c3_598|2>>, a yield of 2{,}800 h gives <<GAP:se_OB_csp_y2800_none:c3s_se_OB_csp_y2800_598|2>>, 3{,}300 h gives <<GAP:se_OB_csp_y3300_none:c3s_se_OB_csp_y3300_598|2>> and an 8-hour capacity credit <<GAP:se_OB_csp_ch8_none:c3s_se_OB_csp_ch8_598|2>> (Table~\ref{tab:sens}). With nuclear, CSP enters only with the 3{,}300 h yield, and then it lowers optimal nuclear from <<RUN:c3_OB_670_h6:N_GW|1>> to <<RUN:c3s_se_OB_csp_y3300_598:N_GW|1>>\,GW (six hours): at fixN 670 nuclear is more sensitive to the CSP yield than at 598. In the other three cases nuclear changes by at most <<EXPR:max(abs(V('RUN:c3s_se_OB_csp_'+k+'_598:N_GW')-V('RUN:c3_OB_670_h6:N_GW')) for k in ('c3','y2800','ch8'))|2>>\,GW.

\paragraph{Imports at double price.} Doubling the price of imports in the critical weeks, instead of the 17\% mark-up, gives gaps of <<GAP:se_PI_pm2_none:c3s_se_PI_pm2_598|2>> (PI) and <<GAP:se_OB_pm2_none:c3s_se_OB_pm2_598|2>>\,EUR/MWh (OB), against <<GAP:val_PI_none:c3_PI_670_h6|2>> and <<GAP:val_OB_none:c3_OB_670_h6|2>>, and optimal nuclear of <<RUN:c3s_se_PI_pm2_598:N_GW|1>>\,GW (PI) and <<RUN:c3s_se_OB_pm2_598:N_GW|1>> (OB): the value of nuclear is insensitive to that price.

\begin{table}[htbp]
\centering
\caption{Sensitivities, six hours, four years, imports with a capacity credit (earlier rule); cost in EUR/MWh.}\label{tab:sens}
\input{tab_sens}
\end{table}

\paragraph{Small modular reactors.} With large nuclear switched off, SMRs are not built at contract costs: Darlington's 17{,}200 and 11{,}700\,EUR/kW correspond to fixN 1{,}179 and 850, above the thresholds. They enter only at two values that are not contracts, 6{,}250 and 3{,}500\,EUR/kW, with 27 and 38\,GW and system costs of 74.78 and 59.20\,EUR/MWh in PI (Table~\ref{tab:smr}). These are developers' and press figures, not offers: boundary scenarios, not results.

\begin{table}[htbp]
\centering
\caption{Small modular reactors with large nuclear off, six hours, four years, imports with a capacity credit. SMR fixN with a 5-year build, 60-year life, $r=5\%$.}\label{tab:smr}
\input{tab_smr}
\end{table}

\paragraph{Net-zero emissions.} In the NZ mode gas with CCS and direct air capture (DAC) are allowed. With nuclear at fixN 670 gas is not chosen in the cases tried (first-plant costs). Without nuclear, gas with CCS does not enter at first-plant costs; at target costs it enters in PI with 6.5\,GW (cost 88.68). Gas without capture with all emissions offset by DAC lowers the cost without nuclear to 85.00 in PI (13.4\,GW of gas, 27.6\,Mt CO$_2$ a year of DAC at target cost) and 83.14 in OB (10.3\,GW, 19.1\,Mt). With target DAC the gap falls from <<GAP:val_PI_none:c3_PI_670_h6|2>> to <<GAP:se_NZ_PI_none_noccs_dactarget:c3s_se_NZ_PI_nuc598|2>> in PI and from <<GAP:val_OB_none:c3_OB_670_h6|2>> to <<GAP:se_NZ_OB_none_noccs:c3s_se_NZ_OB_nuc598|2>> in OB (six hours, against ZD nuclear at fixN 670). The result depends on a DAC cost of 100\,EUR/t (target) or 650 (first plant), our assumptions; 27.6\,Mt a year is about 770 times the nameplate of the largest plant in the world, 36{,}000\,t a year \citep{Climeworks2026}, so the case marks the limit of the assumption favourable to gas rather than an estimate.

\begin{table}[htbp]
\centering
\caption{NZ variants, six hours, four years, imports with a capacity credit.}\label{tab:nz}
\input{tab_nz}
\end{table}

\subsection{Worst weather years}
The four worst PECD years are 2015, 1998, 2014 and 1994 (mean of the three worst weeks of the solar and wind index 13.58 to 14.66\,GW against 16.3 to 19.1 for 2016--2019). Under net self-sufficiency, with them the cost without nuclear rises by <<EXPR:V('RUN:ss0_wx4_PI_none_worst:cost_new')-V('RUN:ss0_PI_none_h6:cost_new')|2>>\,EUR/MWh (PI) and <<EXPR:V('RUN:ss0_wx4_OB_none_worst:cost_new')-V('RUN:ss0_OB_none_h6:cost_new')|2>> (OB), that with nuclear by <<EXPR:V('RUN:ss0_wx4_PI_670_worst:cost_new')-V('RUN:ss0_PI_670_h6:cost_new')|2>> and <<EXPR:V('RUN:ss0_wx4_OB_670_worst:cost_new')-V('RUN:ss0_OB_670_h6:cost_new')|2>>, and optimal nuclear by <<EXPR:V('RUN:ss0_wx4_PI_670_worst:N_GW')-V('RUN:ss0_PI_670_h6:N_GW')|1>> and <<EXPR:V('RUN:ss0_wx4_OB_670_worst:N_GW')-V('RUN:ss0_OB_670_h6:N_GW')|1>>\,GW, to <<RUN:ss0_wx4_PI_670_worst:N_GW|0>>\,GW (<<RUN:ss0_wx4_PI_670_worst:sh|0>>\%) in PI and <<RUN:ss0_wx4_OB_670_worst:N_GW|0>> (<<RUN:ss0_wx4_OB_670_worst:sh|0>>\%) in OB. The gap grows to <<EXPR:V('RUN:ss0_wx4_PI_none_worst:cost_new')-V('RUN:ss0_wx4_PI_670_worst:cost_new')|2>> against <<GAP:ss0_PI_none_h6:ss0_PI_670_h6|2>> in PI and <<EXPR:V('RUN:ss0_wx4_OB_none_worst:cost_new')-V('RUN:ss0_wx4_OB_670_worst:cost_new')|2>> against <<GAP:ss0_OB_none_h6:ss0_OB_670_h6|2>> in OB (Table~\ref{tab:wx}). Nuclear protects against low-wind low-sun periods, which are less severe in the four base years; without help from abroad the system without nuclear covers them with storage and capacity used a few weeks a year.

\begin{table}[htbp]
\centering
\caption{Worst weather years against the base years, six hours, net self-sufficiency, nuclear at fixN 670; cost in EUR/MWh.}\label{tab:wx}
\input{tab_wx}
\end{table}
\FloatBarrier
